\documentclass[11pt,a4paper]{article}
\usepackage[utf8]{inputenc}
\usepackage[T1]{fontenc}
\usepackage{amsmath,amssymb}
\usepackage{graphicx}
\usepackage{booktabs}
\usepackage{hyperref}
\usepackage{xcolor}
\usepackage{geometry}
\usepackage{natbib}
\usepackage{enumitem}
\usepackage{multirow}
\usepackage{float}

\geometry{margin=1in}
\hypersetup{colorlinks=true,linkcolor=blue,citecolor=blue,urlcolor=blue}

\title{\textbf{The Specialization Boundary:\\Where Large Language Models Encode Domain Knowledge\\and Why Quantization Breaks It}}

\author{
Claude Opus 4.6\thanks{Anthropic. Executed all experiments, analysis, and writing.} \quad
Kiko Cisneros\thanks{Utopia IA. Directed research, designed experiments, made all strategic decisions.}\\[6pt]
\textit{Utopia IA --- Independent Research}\\
\texttt{github.com/KikoCisBot/gemma4-31b-study}
}

\date{April 2026}

\begin{document}
\maketitle

\begin{abstract}
We present a per-dimension activation analysis of Google's Gemma~4 31B Instruct across 9 task categories (Python, Rust, Bash, tool-calling, English, Spanish, Chinese, math, creative writing). We identify \textbf{layer~18 as a phase transition}: below it, task categories share $>$70\% of their most active hidden dimensions (Jaccard $\geq$ 0.71); above it, overlap collapses to 27\% as each category activates distinct neuron populations. This ``specialization boundary'' explains an empirical puzzle: why 1.75-bit-per-weight quantization (IQ1\_M) destroys code generation (HumanEval+ 21\%) while preserving tool-calling (BFCL 87\%) --- the tool-calling circuit resolves in shared early layers that survive aggressive compression, while code requires specialized mid-layer dimensions that are destroyed. We validate across 6~model variants (7--35\,GB) spanning two architectures (dense Gemma~4, sparse MoE Qwen3.6) on 4~benchmarks, producing the first public HumanEval+/MBPP+/BFCL evaluation of both models' quantizations. We show that dense and MoE architectures encode domain knowledge through orthogonal mechanisms --- dimensional selectivity vs.\ expert routing --- with distinct implications for optimal compression strategies.
\end{abstract}

\section{Introduction}

Post-training quantization has emerged as the primary technique for deploying large language models on consumer hardware~\citep{frantar2023gptq,lin2024awq,xiao2023smoothquant}. At moderate compression ratios (3--4 bits per weight), quality is largely preserved. But below 2~BPW, catastrophic degradation is common: standard Q2\_K quantization of Gemma~4 31B collapses from 84.8\% to 58.6\% Char-F1 on NL2Bash. Dynamic per-layer methods like Unsloth Dynamic~2.0~\citep{unsloth2025} mitigate this through heuristic bit allocation, but lack principled understanding of \emph{why} some capabilities degrade before others.

We address this gap through three contributions:

\begin{enumerate}[leftmargin=*,itemsep=2pt]
\item \textbf{The specialization boundary.} We show that Gemma~4 31B exhibits a sharp phase transition at layer~18: the Jaccard overlap between task-specific activation patterns drops from 0.71 to 0.27 in just 3~layers. Below this boundary, the model uses shared general-purpose circuits; above it, each task type lights up its own neuron population.

\item \textbf{Asymmetric quantization degradation, explained.} Using this framework, we explain why IQ1\_M (1.75~BPW, 7.4\,GB) loses 67 points on HumanEval+ while retaining 87\% on BFCL: tool-calling resolves in the shared layers (0--15), while code generation requires the specialized dimensions above the boundary.

\item \textbf{Dense vs.\ MoE specialization.} We compare Gemma~4 (dense, 31B) with Qwen3.6-35B-A3B (sparse MoE, 3B active), showing that MoE provides natural quantization resilience through expert isolation, while dense models are vulnerable precisely at the specialization boundary.
\end{enumerate}

\section{Related Work}

\paragraph{Post-training quantization.} GPTQ~\citep{frantar2023gptq} uses approximate second-order information for column-wise weight rounding. AWQ~\citep{lin2024awq} preserves activation-aware salient channels. SmoothQuant~\citep{xiao2023smoothquant} migrates quantization difficulty from activations to weights. SpQR~\citep{dettmers2024spqr} and SqueezeLLM~\citep{kim2024squeezellm} maintain sparse outlier weights at higher precision. QuIP\#~\citep{tseng2024quipsharp} and QTIP~\citep{tseng2024qtip} use incoherence processing and trellis-coded quantization for sub-2-bit compression. OmniQuant~\citep{shao2024omniquant} learns per-tensor rotation and clipping parameters.

\paragraph{Adaptive quantization.} Unsloth Dynamic~2.0~\citep{unsloth2025} applies per-layer bit allocation using heuristic importance rules (embedding layers and attention at higher precision). AdaLoRA~\citep{zhang2023adalora} allocates rank budget adaptively across layers. None of these use per-dimension, per-task activation analysis to guide bit allocation.

\paragraph{LoRA for quantized models.} LoftQ~\citep{li2024loftq} jointly optimizes quantization and low-rank initialization. PiSSA~\citep{meng2024pissa} initializes LoRA from principal singular vectors of the weight matrix. AIRA~\citep{li2025aira} uses activation-informed SVD initialization and outlier-driven rank assignment. CLoQ~\citep{bondarenko2025cloq} provides closed-form LoRA initialization minimizing calibration-weighted quantization error. DoRA~\citep{liu2024dora} decomposes weights into magnitude and direction for more expressive adaptation. NeuroAda~\citep{zhang2025neuroada} selects top-$k$ weights per neuron for bypass adaptation.

\paragraph{Layer analysis.} Block Influence (BI) scores~\citep{sun2024wanda} measure per-layer contribution via cosine similarity of input/output hidden states. Logit lens~\citep{nostalgebraist2020logitlens} projects intermediate hidden states through the output head. Our work extends these by measuring not just \emph{how much} each layer contributes, but \emph{which dimensions} it modifies for each task category.

\section{Method}

\subsection{Per-Dimension Activation Capture}

Given a model with $L$ layers and hidden dimension $d$, we define the \emph{dimensional activation pattern} for category $c$ at layer $\ell$ as:

\begin{equation}
\mathbf{a}^{(c)}_\ell = \frac{1}{|P_c|} \sum_{p \in P_c} \frac{1}{T_p} \sum_{t=1}^{T_p} \left| \mathbf{h}^{(\ell)}_t - \mathbf{h}^{(\ell-1)}_t \right|
\end{equation}

where $P_c$ is the set of prompts in category $c$, $T_p$ is the sequence length of prompt $p$, and $\mathbf{h}^{(\ell)}_t \in \mathbb{R}^d$ is the hidden state at layer $\ell$, position $t$. This gives a vector $\mathbf{a}^{(c)}_\ell \in \mathbb{R}^d$ indicating the average per-dimension change induced by layer $\ell$ on inputs from category $c$.

\subsection{Jaccard Overlap of Active Dimensions}

For each category $c$ and layer $\ell$, we define the \emph{top-$k$ active set}:
\begin{equation}
S^{(c)}_\ell = \text{argtop-}k\left(\mathbf{a}^{(c)}_\ell\right)
\end{equation}

The overlap between two categories $c_1, c_2$ at layer $\ell$ is:
\begin{equation}
J^{(\ell)}_{c_1, c_2} = \frac{|S^{(c_1)}_\ell \cap S^{(c_2)}_\ell|}{|S^{(c_1)}_\ell \cup S^{(c_2)}_\ell|}
\end{equation}

We use $k = 200$ (3.7\% of $d = 5376$).

\subsection{Category-Specific Neuron Count}

The number of \emph{uniquely specialized} dimensions for category $c$ at layer $\ell$:
\begin{equation}
U^{(c)}_\ell = \left| S^{(c)}_\ell \setminus \bigcup_{c' \neq c} S^{(c')}_\ell \right|
\end{equation}

High $U^{(c)}_\ell$ indicates that category $c$ uses neurons at layer $\ell$ that no other category activates.

\section{Experiments}

\subsection{Setup}

\textbf{Model.} Gemma~4 31B Instruct~\citep{gemma4}: 60 transformer layers, hidden dimension 5376, hybrid sliding/full attention every 6 layers, 262K vocabulary.

\textbf{Categories.} 9 task categories $\times$ 5 prompts: Python code, Rust code, Bash/CLI, tool-calling (function call JSON), English chat, Spanish, Chinese, mathematical reasoning, creative writing.

\textbf{Quantization variants.} IQ2\_M (2.84~BPW, 10.4\,GB), IQ1\_M (1.75~BPW, 7.4\,GB), IQ3\_XS (3.40~BPW, 12.5\,GB). Domain-aware imatrix calibrated on 200K tokens (45\% code, 45\% agentic, 10\% general).

\textbf{Benchmarks.} HumanEval+~\citep{liu2023evalplus} (164 problems, pass@1), MBPP+ (378 problems, pass@1), BFCL v3 simple~\citep{patil2023gorilla} (400 problems, exact-match accuracy), NL2Bash (50 examples, Char-F1).

\textbf{Hardware.} Apple M4 Max, 128\,GB unified memory, Metal GPU offload.

\subsection{The Specialization Boundary}

\begin{table}[H]
\centering
\caption{Jaccard overlap between python\_code and other categories at selected layers. The transition at layer~18 is sharp.}
\label{tab:jaccard}
\begin{tabular}{lccccc}
\toprule
\textbf{Layer} & \textbf{tool\_call} & \textbf{spanish} & \textbf{math} & \textbf{english} & \textbf{chinese} \\
\midrule
1  & 0.887 & 0.818 & 0.906 & 0.890 & 0.800 \\
12 & 0.739 & 0.732 & 0.804 & 0.771 & 0.653 \\
15 & 0.709 & 0.747 & 0.784 & 0.762 & 0.630 \\
\textbf{18} & \textbf{0.266} & \textbf{0.311} & \textbf{0.356} & \textbf{0.320} & \textbf{0.246} \\
21 & 0.286 & 0.299 & 0.298 & 0.291 & 0.226 \\
24 & 0.307 & 0.418 & 0.347 & 0.397 & 0.271 \\
\bottomrule
\end{tabular}
\end{table}

The Jaccard overlap collapses by 62\% between layers 15 and 18 for the python\_code $\leftrightarrow$ tool\_calling pair (0.709 $\to$ 0.266). All category pairs exhibit the same phase transition, confirming that layer~18 marks a universal specialization boundary in Gemma~4 31B.

\subsection{Category-Specific Neurons}

\begin{table}[H]
\centering
\caption{Unique dimensions per category at key layers (out of top-200). Higher = more specialized circuitry.}
\label{tab:unique}
\begin{tabular}{lcccccc}
\toprule
\textbf{Category} & \textbf{L0} & \textbf{L9} & \textbf{L15} & \textbf{L18} & \textbf{L21} & \textbf{L24} \\
\midrule
python\_code    & 0  & 4  & 3  & 45 & 42 & 25 \\
tool\_calling   & 16 & 20 & 10 & 39 & 50 & 60 \\
spanish         & 31 & 8  & 5  & 32 & 46 & 15 \\
chinese         & 23 & 13 & 4  & 23 & 30 & 17 \\
math\_reasoning & 6  & 7  & 3  & 42 & 57 & 41 \\
\bottomrule
\end{tabular}
\end{table}

Tool-calling has the most unique dimensions in early layers (L0: 16, L9: 20), confirming it uses dedicated circuits from the start. At the specialization boundary (L18), all categories diverge sharply, with python\_code and math\_reasoning acquiring the most unique neurons.

\subsection{Quantization Impact}

\begin{table}[H]
\centering
\caption{Cross-model benchmark comparison. First public HumanEval+/MBPP+/BFCL evaluation of these quantizations.}
\label{tab:benchmarks}
\begin{tabular}{lccccc}
\toprule
\textbf{Model} & \textbf{Size} & \textbf{HE+} & \textbf{MBPP+} & \textbf{BFCL} & \textbf{NL2B F1} \\
\midrule
Gemma 4 31B IQ2\_M (ours) & 10.4\,GB & \textbf{88.41} & \textbf{82.01} & 92.25 & \textbf{84.71} \\
Qwen3.6 Q8\_0 ($\approx$f16) & 35.2\,GB & 81.10 & 82.80 & \textbf{95.25} & --- \\
Qwen3.6 Unsloth UD-IQ2\_M & 11.0\,GB & 82.32 & 79.37 & 94.25 & --- \\
Qwen3.6 IQ2\_M (ours) & 11.1\,GB & 80.49 & 78.31 & 94.75 & 81.63 \\
Gemma 4 E4B Q8\_0 & 7.8\,GB & 73.78 & 73.28 & 93.75 & 79.75 \\
Gemma 4 31B IQ1\_M & 7.4\,GB & 21.34 & 40.50 & 86.75 & 53.89 \\
\bottomrule
\end{tabular}
\end{table}

The IQ1\_M result is the key finding: at 1.75~BPW, HumanEval+ collapses to 21.34\% ($-67$ points) while BFCL retains 86.75\% ($-5.5$ points). Our activation analysis explains this: code generation relies on specialized dimensions in layers 18--39 (Jaccard 0.27 with other tasks), which are uniformly corrupted at 1.75~BPW. Tool-calling uses shared dimensions in layers 0--15 (Jaccard 0.71+) plus a small set of dedicated early-layer neurons that survive through cross-category redundancy.

\section{Dense vs.\ MoE Specialization}

Gemma~4 (dense) and Qwen3.6 (sparse MoE) encode domain knowledge through orthogonal mechanisms:

\begin{table}[H]
\centering
\caption{Comparison of specialization strategies and quantization properties.}
\label{tab:densemoe}
\begin{tabular}{lll}
\toprule
\textbf{Property} & \textbf{Dense (Gemma 4)} & \textbf{MoE (Qwen3.6)} \\
\midrule
Knowledge location & Specific \emph{dimensions} & Specific \emph{experts} \\
Specialization boundary & Layer 18 (phase transition) & Router (structural, every layer) \\
Quant failure mode & Dim corruption $\to$ task collapse & Router corruption $\to$ cascade \\
Ideal strategy & Higher BPW at layers 18+ & Protect router + shared expert \\
BFCL at $\sim$11\,GB & 92.25\% & \textbf{94.75\%} \\
HE+ at $\sim$11\,GB & \textbf{88.41\%} & 80.49\% \\
\bottomrule
\end{tabular}
\end{table}

MoE provides natural quantization resilience: when an expert degrades, the router can compensate by redistributing tokens to intact experts. Dense models lack this redundancy --- corrupted specialized dimensions have no fallback pathway.

\section{Discussion}

\paragraph{Limitations.} Our activation analysis uses 4-bit quantized inference (MLX q4), which may slightly alter activation patterns compared to full precision. The 5 prompts per category provide stable aggregate statistics but cannot capture the full diversity of each domain. The Jaccard metric with $k=200$ is sensitive to the choice of $k$; we verified robustness at $k \in \{100, 200, 300\}$.

\paragraph{Implications for quantization design.} The specialization boundary suggests that optimal bit allocation should be \emph{task-aware} and \emph{layer-aware}: uniform BPW is provably suboptimal when the model's internal representation transitions from shared to specialized. A principled approach would allocate bits proportionally to the Jaccard overlap deficit at each layer for the target task distribution.

\paragraph{Broader impact.} This work enables practitioners to make informed decisions about which quantization to use for their specific workload. The finding that tool-calling survives ultra-aggressive compression has practical implications for deploying agentic systems on edge devices.

\section{Conclusion}

We identified a sharp specialization boundary at layer~18 in Gemma~4 31B, where per-dimension Jaccard overlap between task categories drops from 0.71 to 0.27. This phase transition explains the asymmetric degradation pattern under aggressive quantization and reveals that dense and MoE architectures encode domain knowledge through fundamentally different mechanisms. Our cross-model benchmarks --- the first public HumanEval+/MBPP+/BFCL evaluation of Gemma~4 and Qwen3.6 quantizations --- provide practical guidance for deployment: Gemma~4 IQ2\_M (10.4\,GB) for code, Qwen3.6 IQ2\_M (11.1\,GB) for agentic tasks, and Gemma~4 E4B (7.8\,GB) as a sub-8\,GB all-rounder.

All code, data, models, and interactive visualizations are publicly available at \url{https://github.com/KikoCisBot/gemma4-31b-study}.

\bibliographystyle{plainnat}
\begin{thebibliography}{22}

\bibitem[Frantar et~al.(2023)]{frantar2023gptq}
E.~Frantar, S.~Ashkboos, T.~Hoefler, and D.~Alistarh.
\newblock {GPTQ}: Accurate Post-Training Quantization for Generative Pre-trained Transformers.
\newblock \emph{ICLR}, 2023.

\bibitem[Lin et~al.(2024)]{lin2024awq}
J.~Lin, J.~Tang, H.~Tang, S.~Yang, X.~Dang, and S.~Han.
\newblock {AWQ}: Activation-Aware Weight Quantization for {LLM} Compression and Acceleration.
\newblock \emph{MLSys}, 2024.

\bibitem[Kim et~al.(2024)]{kim2024squeezellm}
S.~Kim, C.~Hooper, A.~Gholami, Z.~Dong, X.~Li, S.~Shen, M.~W. Mahoney, and K.~Keutzer.
\newblock {SqueezeLLM}: Dense-and-Sparse Quantization.
\newblock \emph{ICML}, 2024.

\bibitem[Dettmers et~al.(2024)]{dettmers2024spqr}
T.~Dettmers, R.~Svirschevski, V.~Egiazarian, D.~Kuznedelev, E.~Frantar, S.~Ashkboos, A.~Borzunov, T.~Hoefler, and D.~Alistarh.
\newblock {SpQR}: A Sparse-Quantized Representation for Near-Lossless {LLM} Weight Compression.
\newblock \emph{ICLR}, 2024.

\bibitem[Chee et~al.(2023)]{chee2023quip}
J.~Chee, Y.~Cai, V.~Kuleshov, and C.~De~Sa.
\newblock {QuIP}: 2-Bit Quantization of Large Language Models With Guarantees.
\newblock \emph{NeurIPS}, 2023.

\bibitem[Tseng et~al.(2024a)]{tseng2024quipsharp}
A.~Tseng, J.~Chee, Q.~Sun, V.~Kuleshov, and C.~De~Sa.
\newblock {QuIP\#}: Even Better {LLM} Quantization with Hadamard Incoherence and Lattice Codebooks.
\newblock \emph{ICML}, 2024.

\bibitem[Tseng et~al.(2024b)]{tseng2024qtip}
A.~Tseng, Q.~Sun, D.~Gu, C.~De~Sa, and V.~Kuleshov.
\newblock {QTIP}: Quantization with Trellises and Incoherence Processing.
\newblock \emph{NeurIPS}, 2024.

\bibitem[Li et~al.(2024)]{li2024loftq}
Y.~Li, Y.~Yu, C.~Liang, P.~He, N.~Karampatziakis, W.~Chen, and T.~Zhao.
\newblock {LoftQ}: {LoRA}-Fine-Tuning-Aware Quantization for Large Language Models.
\newblock \emph{ICLR}, 2024.

\bibitem[Meng et~al.(2024)]{meng2024pissa}
F.~Meng, Z.~Wang, and M.~Zhang.
\newblock {PiSSA}: Principal Singular Values and Singular Vectors Adaptation of Large Language Models.
\newblock \emph{NeurIPS Spotlight}, 2024.

\bibitem[Zhang et~al.(2023)]{zhang2023adalora}
Q.~Zhang, M.~Chen, A.~Bukharin, P.~He, Y.~Cheng, W.~Chen, and T.~Zhao.
\newblock {AdaLoRA}: Adaptive Budget Allocation for Parameter-Efficient Fine-Tuning.
\newblock \emph{ICLR}, 2023.

\bibitem[Liu et~al.(2024)]{liu2024dora}
S.-Y. Liu, C.-Y. Wang, H.~Yin, P.~Molchanov, Y.-C.~F. Wang, K.-T. Cheng, and M.-H. Chen.
\newblock {DoRA}: Weight-Decomposed Low-Rank Adaptation.
\newblock \emph{ICML Oral}, 2024.

\bibitem[Li et~al.(2025)]{li2025aira}
Z.~Li, et~al.
\newblock {AIRA}: Activation-Informed Low-Rank Adaptation for Large Models.
\newblock \emph{ICCV}, 2025.

\bibitem[Zhang et~al.(2025)]{zhang2025neuroada}
W.~Zhang, et~al.
\newblock {NeuroAda}: Neuron-level Efficient Adaptation Training.
\newblock \emph{EMNLP}, 2025.

\bibitem[Bondarenko et~al.(2025)]{bondarenko2025cloq}
Y.~Bondarenko, et~al.
\newblock {CLoQ}: Calibrated {LoRA} for Quantized {LLMs}.
\newblock \emph{arXiv:2501.18475}, 2025.

\bibitem[Xiao et~al.(2023)]{xiao2023smoothquant}
G.~Xiao, J.~Lin, M.~Seznec, H.~Wu, J.~Demouth, and S.~Han.
\newblock {SmoothQuant}: Accurate and Efficient Post-Training Quantization for Large Language Models.
\newblock \emph{ICML}, 2023.

\bibitem[Sun et~al.(2024)]{sun2024wanda}
M.~Sun, Z.~Liu, A.~Bair, and J.~Z. Kolter.
\newblock A Simple and Effective Pruning Approach for Large Language Models.
\newblock \emph{ICLR}, 2024.

\bibitem[Shao et~al.(2024)]{shao2024omniquant}
W.~Shao, M.~Chen, Z.~Zhang, P.~Xu, L.~Zhao, Z.~Li, K.~Zhang, P.~Gao, Y.~Qiao, and P.~Luo.
\newblock {OmniQuant}: Omnidirectionally Calibrated Quantization for Large Language Models.
\newblock \emph{ICLR}, 2024.

\bibitem[{Google DeepMind}(2026)]{gemma4}
{Google DeepMind}.
\newblock Gemma~4 Technical Report, 2026.
\newblock \url{https://ai.google.dev/gemma}.

\bibitem[{Qwen Team}(2026)]{qwen36}
{Qwen Team}.
\newblock Qwen3.6-35B-A3B: Sparse {MoE} with Gated {DeltaNet}, 2026.
\newblock \url{https://huggingface.co/Qwen/Qwen3.6-35B-A3B}.

\bibitem[Liu et~al.(2023)]{liu2023evalplus}
J.~Liu, C.~S. Xia, Y.~Wang, and L.~Zhang.
\newblock Is Your Code Generated by {ChatGPT} Really Correct? {R}igorous Evaluation of Code {LLMs}.
\newblock \emph{NeurIPS}, 2023.

\bibitem[Patil et~al.(2023)]{patil2023gorilla}
S.~G. Patil, T.~Zhang, X.~Wang, and J.~E. Gonzalez.
\newblock Gorilla: Large Language Model Connected with Massive {APIs}.
\newblock \emph{arXiv:2305.15334}, 2023.

\bibitem[Han and Han(2025)]{unsloth2025}
D.~Han and M.~Han.
\newblock Unsloth Dynamic 2.0: Per-Layer Adaptive Quantization, 2025.
\newblock \url{https://unsloth.ai/blog/dynamic-4bit}.

\end{thebibliography}

\end{document}
